\documentstyle[% 


righttag% 


,11pt% 


,amssymb% 


,russian%               remove in Eng. 


,izvru%                 Set izven% in Eng. 


% ,izvru-ks             for brief communications 


]{amsart} 


 


%       Math definitions 


 


\newcommand{\la}{\langle} 


\newcommand{\ra}{\rangle} 


\newcommand{\tts}[1]{} 


\renewcommand{\baselinestretch}{1.05}


 


%\hoffset=-15mm%                  For Eng. 


 


\setcounter{page}{13} 


\god{2000} 


\nomer{2~(453)} %                        Remove in Eng. 


%\nomer{Vol.\,44, No.\,2}%               Set in Eng. 


%\str{\pageref{begin}--\pageref{end}}%  Set in Eng. 


\udk{517.984 } 


 


\author{\it А.М.\,Ахтямов} 


% NB:   ^^ remove in Eng. 


\title{О коэффициентах разложений по собственным функциям  


краевой задачи со спектральным параметром  


в граничных условиях} 


 


\date{} 


% \pagestyle{myheadings}%         Set in Eng. 


\pagestyle{plain} %              Remove in Eng. 


\begin{document} 


% \thispagestyle{plain}%          Set in Eng. 


\maketitle 


\label{begin} 


 


\section{Введение} 


 


Задача определения силы тока $i(x,t)$ в проводе, концы  


которого заземлены через сосредоточенные сопротивления, имеет  


вид 


\begin{gather*} 


i_{xx}=C L i_{tt}+(C R+G L) i_t+ G R i,\quad 


i_x(0,t)-C R^{(1)}_0 i_t(0,t)- 


G R^{(1)}_0 i(0,t)=0,\\ 


% 


i_x(l,t)-C R^{(2)}_0 i_t(l,t)- 


G R^{(2)}_0 i(l,t)=0,\quad i(x,0)=\varphi(x),\\ 


% 


i_x(x,0)=(-R\varphi(x)-f'(x))/L 


\end{gather*} 


([1], cc.\,21, 176). Здесь $\varphi(x)$ и $f(x)$ --- 


соответственно сила тока и напряжение в  


начальный момент времени.  


 


Рассматриваемая задача допускает разделение переменных по  


времени, поэтому для ее решения пригоден метод Фурье. При  


этом у соответствующей задачи в краевых условиях  


появляется параметр, в результате чего нарушается минимальность  


цепочек собственных функций в пространстве $L_2\times L_2$ ([2], 


c.\,191). 


Обоснование разложимости в ряды по цепочкам собственных функций  


для таких задач появилось сравнительно недавно ([2], c.\,190--229). 


Проводится оно в пространствах более сложной структуры.  


Задача вычисления коэффициентов соответствующих рядов не  


является тривиальной. Коэффициенты подобных рядов  


находились лишь для некоторых спектральных задач ([3], c.\,137--139) 


с помощью сопряженного оператора. 


 


Вместо указанной конкретной задачи ниже рассмотрим общую  


краевую задачу второго порядка и для нее найдем коэффициенты  


разложения пары произвольных функций (принадлежащих специальным  


пространствам [2]) в ряды по собственным функциям. Cделаем это  


двумя способами: первый подход обобщает [3] и использует явное  


построение сопряженного оператора к линеаризатору А.А.\,Шкаликова  


задачи (1)--(3), а второй подход использует возможность  


представления задачи в виде пучка неограниченных операторов с 


дальнейшим применением соотношений биортогональности [4],


известных для пучков операторов.  


 


Для удобства приведем список обозначений. 


Через $\Bbb C$ обозначается пространство комплексных чисел; 


$L_2$ --- пространство суммируемых с квадратом функций,  


заданных на отрезке $[0,1]$; $W^1_2$ --- соболевское пространство  


функций на отрезке $[0,1]$, $\cal W^0_2$ --- пространство $W^1_2\times 


L_2$, 


$\frak H$ --- пространство $(L_2\times\Bbb C^2)^2$. Через $(\ ,\ )$, 


$[\ ,\ ]$, $\la\ ,\ \ra$, $(\ ,\ )_{\frak H}$ 


записываются скалярные произведения в пространствах $L_2$, 


$L_2\times\Bbb C^2$, $\cal W^0_2$, $\frak H$ 


соответственно, а $y$, $\bold y$, $\wt{\bold v}$, $\wh{\bold y}$ --- 


соответственно элементы этих пространств.


 


\section{Вычисление коэффициентов с помощью сопряженного оператора} 


 


Рассмотрим спектральную задачу  


\begin{align} 


l(y,\lambda)&=y''+(p_0+p_1\lambda)y'+ 


(q_0+q_1\lambda+q_2\lambda^2)y=0,\\ 


% 


U_1(y,\lambda)&=a_1y'(0)+a_2y(0)+\lambda y(0)=0, \\ 


% 


U_2(y,\lambda)&=b_1y'(1)+b_2y(1)+\lambda y(1)=0. 


\end{align} 


Здесь $a_1,a_2,b_1,b_2,q_2\in\Bbb C$; $a_1,b_1,q_2\ne0$; $p_0=p_0(x)$, 


$q_0=q_0(x)$, $p_1=p_1(x)$, $q_1=q_1(x)$ --- 


непрерывно дифференцируемые комплекснозначные функции.  


 


Сопряженной к задаче (1)--(3) (см., напр., [5], c.\,17--23) является  


краевая задача 


\begin{align*} 


z''+(\ov p_0z)'+\ov q_0z+((-\ov p_1z)'+\ov q_1)\mu^2+ 


\ov q_2\mu z&=0,\\ 


% 


\ov a_1z'(0)+(\ov a_2-\ov a_1\ov{p_0(0)})z(0)+ 


(1-\ov a_1\ov{p_0(0)})\mu z(0)&=0,\\ 


% 


\ov b_1z'(1)+(\ov b_2-\ov b_1\ov{p_0(1)})z(1)+ 


(1-\ov b_1\ov{p_1(1)})\mu z(1)&=0. 


\end{align*} 


 


Согласно [2] задача (1)--(3) допускает линеаризацию в  


пространстве $\cal W^0_2=W^1_2\times L_2$. Линеаризатор $\cal H$ задачи 


(1)--(3) 


задается равенствами 


\begin{gather*} 


\cal H\{v_0,v_1\}=\{v_1,\ -q^{-1}_2(v_0''+p_0v'_0+ 


q_0v_0+p_1v'_1+q_1v_1)\}, \\


% 


D(\cal H)=\Big\{\wt{\bold v}=\{v_0,v_1\}\mid v_0\in W^2_2, 


\ v_1\in W^1_2, \\ 


% 


V_1(\wt{\bold v})=a_1v'_2(0)+a_2v_0(0)+v_1(0)=0,\ \;


V_2(\wt{\bold v})=b_1v'_0(1)+ 


b_2v_0(1)+v_1(1)=0\Big\}. 


\end{gather*} 


Линейная спектральная задача $\cal H\wt{\bold v}=\lambda\wt{\bold v}$ 


имеет те же собственные  


значения, что и задача (1)--(3), а собственные функции 


оператора $\cal H$, соответствующие собственным значениям $\lambda_k$, 


имеют вид  


$\wt{\bold v}_k=\{y_k,\lambda_k y_k\}$, где $y_k$ --- собственные 


функции задачи (1)--(3), 


отвечающие тем же собственным значениям $\lambda_k$ (чтобы не  


усложнять существо дела, считаем все собственные значения  


задачи (1)--(3) простыми). 


 


Пусть $\wt{\bold v}_k$ --- собственная функция оператора $\cal H$, 


соответствующая  


собственному значению $\lambda_k$, а $\wt{\bold g}_k$ --- собственная 


функция оператора  


$\cal H^*$. Тогда справедливы соотношения биортогональности 


$\la\wt{\bold v}_k,\wt{\bold g}_j\ra/\la\wt{\bold v}_k, 


\wt{\bold g}_k\ra=\delta_{k,j}$. 


Из этих соотношений получается формула для вычисления  


коэффициентов разложения элемента $\wt{\bold v}=\{v_0,v_1\}$ из 


пространства 


$\cal W^0_2=W^1_2\times L_2$ в ряд по собственным функциям 


$\wt{\bold v}_k$ 


\begin{equation} 


c_k=\la\wt{\bold v},\wt{\bold g}_k\ra/ 


\la\wt{\bold v}_k,\wt{\bold g}_k\ra. 


\end{equation} 


 


Таким образом, задача вычисления коэффициентов сводится к  


поиску сопряженного оператора и его собственных функций. 


 


\begin{thm} 


Сопряженный к $\cal H$ оператор $\cal H^*$ имеет вид 


$\cal H^*\{g_0,g_1\}=H_1\{g_0,g_1\}+H_2\{g_0,g_1\}$, 


\begin{align*} 


D(\cal H^*)&=\Big\{\wt{\bold g}=\{g_0,g_1\}\mid g_0\in W^2_2, 


\ g_1\in W^1_2,\\ 


% 


&\qquad G_1(\wt{\bold g})=\ov q_2^{\,-1}(1-\ov a_1 


\ov p_1(0))g_1(0)-\ov a_1(g_0(0)-g'_0(0))=0,\\ 


% 


&\qquad G_2(\wt{\bold g})=\ov q_2^{\,-1}(1-\ov a_1 


\ov p_1(1))g_1(1)-\ov b_1(g_0(1)+g'_0(1))=0\Big\}. 


\end{align*} 


Здесь  


\begin{align*} 


H_1\{g_0,g_1\}&=\bigg\{\ov q_2^{\,-1}\bigg[g_1- 


\int^x_0\ov{p_0(t)}g_1(t)dt+\int^x_0\int^\xi_0\ov{q_0(\xi)} 


g_1(\xi)d\xi\,dt-\frac{2-x}{3}g_1(0)+\\ 


% 


&\quad+\frac{1+x}{3}\bigg(\int^1_0\ov{p_0(t)}g_1(t)dt- 


\int^1_0\int^t_0\ov{q_0(\xi)}g_1(\xi)d\xi\,dt- 


g_1(1)-\int^1_0\ov{q_0(\xi)}g_1(\xi)d\xi\bigg)\bigg], \\ 


% 


&\qquad -g''_0+\ov q_2^{\,-1}(\ov{p_1(x)}g_1(x))'- 


\ov q_2^{\,-1}(\ov{q_1(x)}g_1(x))\bigg\}, \\ 


% 


H_2\{g_0,g_1\}&=\bigg\{\ov a_2\frac{2-x}{3}(g'_0(0)-g_0(0)- 


\ov q_2^{\,-1}\ov{p_1(0)}g_1(0)) +\\ 


&\quad +\ov b_2\frac{1+x}{3}(-g'_0(1)-g_0(1)+ 


\ov q_2^{\,-1}\ov{p_1(1)}g_1(1)),\ 0\bigg\}. 


\end{align*} 


\end{thm} 


 


{\bf Доказательство.} В отличие от стандартной нормы  


$\|f\|_{W^1_2}=(\|f'\|^2_{L_2}+ \|f\|^2_{L_2})^{1/2}$ 


будем использовать норму  


$\|f\|_{W^1_2}=(\|f'\|^2_{L_2}+|f(0)|^2+|f(1)|^2)^{1/2}$. 


Эти нормы эквивалентны ([6], c.\,147). В нашем случае вторая норма  


удобнее, т.\,к. она позволяет использовать вместо функций Грина  


обычные интегралы с переменным верхним пределом.  


 


Верно равенство 


\begin{equation} 


\la\cal H\wt{\bold v},\wt{\bold g}\ra=\la\wt{\bold v}, 


H_1\wt{\bold g}\ra+P(\wt{\bold v},\wt{\bold g}). 


\end{equation} 


Здесь $\cal H$ и $H_1$ --- выражения, определенные в формулировке  


теоремы, а 


$P(\wt{\bold v},\wt{\bold g})=q^{-1}_2[v'_0(0)\ov g_1(0)- 


v'_0(1)\ov g_1(1)+


v_1(0)(q_2\ov g_0(0)-q_2\ov{g'}_0(0)+p_1(0)\ov g_1(0))+ 


v_1(1)(q_2\ov g_0(1)-q_2\ov{g'}_0(1)+p_1(1)\ov g_1(1))]$


--- билинейная форма. Равенство (5) 


получается интегрированием по частям. Оно  


напоминает формулу Лагранжа, с помощью которой находится  


сопряженный оператор для обыкновенных дифференциальных уравнений  


в пространстве $L_2$. Отличие состоит в том, что в нашем случае в  


``краевых условиях'' $V_1(\wt{\bold v})=0$, $V_2(\wt{\bold v})=0$ 


(точнее, в условиях  


на область определения оператора) содержaтся шесть переменных  


(это переменные $v_0(0)$, $v_0(1)$, $v'_0(0)$, $v'_0(1)$, $v_1(0)$, 


$v_1(1)$), а в  


форме $P(\wt{\bold v},\wt{\bold g})$ --- лишь четыре ($v'_0(0)$, 


$v'_0(1)$, $v_1(0)$, $v_1(1)$). В  


классическом случае --- обратная ситуация. Там количество  


переменных в краевых условиях не больше числа переменных,  


содержащихся в билинейной форме проинтегрированных членов. Поэтому  


методы, используемые в ([5], c.\,17--22),


непосредственно не применимы в нашем случае. Прежде чем  


применить известные методы нахождения сопряженного оператора,  


требуется сделать некоторые преобразования. А именно, форму  


$P(\wt{\bold v},\wt{\bold g})$ 


представим в виде суммы двух форм 


$P_1(\wt{\bold v},\wt{\bold g})$ и $P_2(\wt{\bold v},\wt{\bold g})$, 


первую из которых запишем в виде скалярного произведения


\begin{multline*} 


P_1(\wt{\bold v},\wt{\bold g})=-a_2v_0(0)q^{-1}_2(q_2\ov g_0(0)- 


q_2\ov{g'}_0(0)+p_1(0)\ov g_1(0))- \\ 


% 


-b_2v_0(1)q^{-1}_2(q_2\ov g_0(1)-q_2\ov{g'}_0(1)+ 


p_1(1)\ov g_1(1))=\la\wt{\bold v},H_1\wt{\bold g}\ra, 


\end{multline*} 


а из второй методом, описанным в [5], найдутся  


``сопряженные краевые условия''. 


 


Из изложенного выше следует равенство  


$\la\cal H\wt{\bold v},\wt{\bold g}\ra=\la\wt{\bold v},\cal M\wt{\bold 


g}\ra$, где $\cal M$ --- оператор, который 


определяется равенствами 


$$ 


\cal M\{g_0,g_1\}=H_1\{g_0,g_1\}+H_2\{g_0,g_1\},


$$ 


\begin{align*} 


D(\cal M)&=\Big\{\wt{\bold g}=\{g_0,g_1\}\mid g_0\in W^2_2, 


\ g_1\in W^1_2,\\ 


% 


&\qquad G_1(\wt{\bold g})=\ov q_2^{\,-1}(1-\ov a_1 


\ov p_1(0))g_1(0)-\ov a_1(g_0(0)-g'_0(0))=0,\\ 


% 


&\qquad G_2(\wt{\bold g})=\ov q_2^{\,-1}(1-\ov a_1 


\ov p_1(1))g_1(1)-\ov b_1(g_0(1)+g'_0(1))=0\Big\}. 


\end{align*} 


 


Таким образом, $\cal M\subset\cal H^*$. 


Обратное включение $\cal M\supset\cal H^*$ следует из того факта, что 


образом  


операторов $\cal M$ и $\cal H$ является все пространство $\cal W^0_2$. 


\medskip 


 


{\bf Замечание.} Как видим, формула действия сопряженного  


оператора помимо производных содержит также интегралы и  


функционалы вида $Ax+B$ ($A,B\in\Bbb C$).  


В этом проявляется  


отличие от вида сопряженного оператора к оператору, порожденному 


обыкновенным дифференциальным выражением в пространстве $L_2$.  


 


\begin{thm} 


Коэффициенты разложения произвольной  


функции $\wt{\bold v}$ из пространства $\cal W^0_2$ по элементам 


$\wt{\bold v}_k$ имеют 


вид $c_k=P_k/Q_k$, где  


\begin{align*} 


P_k&=(v_0,\ \ov\lambda_k z_k+\ov q_2^{\,-1}(\ov q_1(x)- 


(\ov p_1(x)z_k(x))'))+


(v_1,z_k)+\lambda^{-1}_k q^{-1}_2 v_0(1)(\ov{z'}_k(1)- 


p_0(1)\ov z_k(1)+ \\ 


% 


&\quad +\tfrac{b_2}{b_1}\ov z_k(1))+\lambda^{-1}_k 


q^{-1}_2 v_0(0)(p_0(0)\ov{z'}_k(0)-\ov{z'}_k(0)- 


\tfrac{a_2}{a_1}\ov z_k(0)), \\ 


% 


Q_k&=(y_k,\ \ov\lambda_k z_k+\ov q_2^{\,-1}(\ov q_1(x)- 


(p_1(x)z_k(x))'))+(\lambda_k y_k,z_k) 


+\lambda^{-1}_k q^{-1}_2 y_k(1)(\ov{z'}_k(1)- 


p_0(1)\ov z_k(1)+ \\ 


% 


&\quad +\tfrac{b_2}{b_1}\ov z_k(1))+\lambda^{-1}_k 


q^{-1}_2 y_k(0)(p_0(0)\ov z_k(0)-\ov{z'}_k(0)- 


\tfrac{a_2}{a_1}\ov z_k(0)).\end{align*} 


\end{thm} 


 


{\bf Доказательство.} С найденным представлением 


сопряженного оператора можно показать, что 1)~собственные значения  


оператора $\cal H^*$ и сопряженной задачи совпадают; 2)~между  


собственными функциями $\wt{\bold g}_k=\{g_{k,0},g_{k,1}\}$ оператора 


$\cal H^*$,  


соответствующими собственным значениям $\mu_k$, и собственными  


функциями $z_k$ сопряженной задачи, соответствующими тем же  


собственным значениям, можно установить взаимно однозначное  


соответствие.  


 


Подставим функции $\wt{\bold g}_k$, найденные из равенства $\cal 


H^*\wt{\bold g}_k=\mu_k\wt{\bold g}_k$, 


в формулу (4) для коэффициентов $c_k$. Обозначим 


$P_k=\la\wt{\bold v},\wt{\bold g}_k\ra$, 


$Q_k=\la\wt{\bold v}_k,\wt{\bold g}_k\ra$. 


Применив формулу интегрирования по частям и  


воспользовавшись тем, что $\mu_k=\ov\lambda_k$, 


получим утверждение теоремы. 


 


\section{Вычисление коэффициентов с помощью сведения  


задачи\protect\\ к квадратичному пучку операторов} 


 


Предложим теперь другой метод вычисления коэффициентов  


разложения, который основан на приведении спектральной задачи  


(1)--(3) к пучку операторов, действующему в пространстве $L_2\times\Bbb 


C^2$. 


 


Пространство $L_2\times\Bbb C^2$ --- гильбертово пространство с  


элементами $\bold y=\{y(x),\xi_1,\xi_2\}$ и скалярным произведением 


$[\bold y,\bold z]=(y(x),z(x))+\xi_1\ov\eta_1+\xi_2\ov\eta_2$. 


Здесь $\bold z=\{z(x),\eta_1,\eta_2\}$. 





Определим в пространстве $L_2\times\Bbb C^2$ операторы $A_0$, $A_1$, 


$A_2$ равенствами 


\begin{align*} 


A_0{\bold y}{} & =\{y''+p_0y'+q_0y,\ a_1y'(0)+a_2y(0), 


\ b_1y'(1)+b_2y(1)\}, \\ 


% 


D(A_0){} & =\{\bold y=\{y(x),\ y(0),\ y(1)\}\mid y\in W^2_2\}, \\ 


% 


A_1{\bold y}{} & = \{p_1y'+q_1y,\ y(0),\ y(1)\},\quad D(A_1)=D(A_0), \\ 


% 


A_2{\bold y}{} & = \{q_2y,\ 0,\ 0\},\quad D(A_2)=D(A_0). 


\end{align*} 


Образуем из этих операторов пучок  


\begin{equation} 


A(\lambda)=A_0+\lambda A_1+\lambda^2A_2. 


\end{equation} 


Собственные значения пучка операторов $A(\lambda)$ и краевой задачи  


совпадают. Между собственными элементами $y_k$ краевой задачи и  


собственными элементами $\bold y_k$ пучка операторов можно установить  


взаимно однозначное соответствие $\bold y_k=\{y_k,\ y_k(0),\ y_k(1)\}$. 


 


Сопряженный к оператору $A_0$ оператор $A^*_0$ имеет вид 


\begin{align*} 


A^*_0\bold z&=\{z''-(p_0z)'+q_0z,\ \ov a_1^{\,-1}\ov a_2z(0)- 


p_0(0)z(0)+z'(0),\ \ov b_1^{\,-1}\ov b_2z(1)-p_0(1)z(1)+ 


z'(1)\},\\ 


% 


D(A^*_0)&=\{\bold z=\{z,\ \ov a_1^{\,-1}z(0), 


\ \ov b_1^{\,-1}z(1)\mid z\in W^2_2\}. 


\end{align*} 


Это следует из равенства $[A_0\bold y,\bold z]=[\bold y,A^*_0\bold z]$ 


и из того, что  


образом операторов $A_0$ и $A_0^*$ является все пространство.  


 


Зададим операторы $M_1$, $M_2$ с помощью равенств 


\begin{align*} 


M_1\{z,d_1,d_2\}&=\{-(\ov p_1z)'+\ov q_1z,\ d_1-\ov p_1(0), 


\ d_2+\ov p_1(1)z(1)\},\\ 


% 


D(M_1)&=\big\{\{z, d_1, d_2\}\mid z\in W^1_2;\ d_1,d_2\in\Bbb C\big\},\\ 


% 


M_2\{z,d_1,d_2\}&=\{q^{-1}_2z,\ 0,\ 0\},\\ 


% 


D(M_2)&=\big\{\{z,d_1,d_2\}\mid z\in L_2;\ d_1,d_2\in\Bbb C\big\}. 


\end{align*} 


Непосредственной проверкой убедимся в том, что 


$[A_1\bold y,\bold z]=[\bold y,M_1\bold z]$, 


$[A_2\bold y,\bold z]=[\bold y,M_2\bold z]$. 


Области определения операторов $M_1$, $M_2$ шире области определения  


оператора $A^*_0$.  


 


Образуем из операторов $A^*_0$, $M_1$, $M_2$ пучок  


\begin{equation} 


M(\mu)=A^*_0+\mu M_1+\mu^2M_2. 


\end{equation} 


Область определения пучка операторов (7) совпадает с областью  


определения оператора $A^*_0$. Заметим, что собственные значения пучка  


операторов (7) и собственные значения сопряженной краевой задачи  


совпадают. Между собственными функциями $\bold z_k$ пучка операторов 


(7),  


соответствующими собственным значениям $\ov\lambda_k$, и собственными  


функциями $z_k$ сопряженной краевой задачи, соответствующими тем же  


$\ov\lambda_k$, можно установить взаимно однозначное соответствие по 


формуле 


$$ 


\bold z_k=\{z_k,\ \ov a_1^{\,-1}z_k(0), 


\ \ov b_1^{\,-1}z_k(1)\}. 


$$ 


Для $\mu\ne\ov\lambda_k$ оператор $M(\mu)$ обратим, поэтому 


$M(\mu)=[A(\lambda)]^*$. 


 


В [4] получены соотношения  


биортогональности для собственных функций пучка неограниченных  


операторов. Для нашего случая эти соотношения имеют вид 


\begin{equation} 


(G\wh{\bold y}_k,\wh{\bold z}_j)_{\frak H}=\delta_{k,j}, 


\end{equation} 


где $\wh{\bold y}_k=\{\bold y_k,\ \lambda_k\bold y_k\}$, 


$\wh{\bold z}_j=\{\bold z_j,\ \mu_j\bold z_j\}$, 


$G=\left(\begin{smallmatrix} A_1 & A_2 \\ A_2 & 


0\end{smallmatrix}\right)$. 


Из (8) следует, что коэффициенты  


разложения произвольной функции $\wh{\bold y}$ из пространства $\frak 


H$ по  


собственным элементам пучка (6) находятся по формуле 


\begin{equation} 


\wh c_k=(G\wh{\bold y},\wh{\bold z}_k)_{\frak H}/ 


(G\wh{\bold y}_k,\wh{\bold z}_k)_{\frak H}. 


\end{equation} 


Если задача (1)--(3) усиленно регулярна (определение см. в  


[2], c.\,196), то любой элемент $\wt{\bold v}=\{v_0,v_1\}$ пространства 


$\cal W^0_2$ 


может быть разложен в ряд по элементам 


$\wt{\bold v}_k=\{y_k,\ \lambda_k y_k\}$ 


\begin{equation} 


\wt{\bold v}=\sum^\infty_{k=1}c_k\wt{\bold v}_k, 


\end{equation} 


который сходится в норме пространства $\cal W^0_2$.  


Отсюда следует  


\begin{equation} 


\wh{\bold y}=\sum^\infty_{k=1}c_k\wh{\bold y}_k, 


\end{equation} 


где $\wh{\bold y}=\{v_0,v_0(0),v_0(1),v_1,v_1(0),v_1(1)\}$, 


$\wh{\bold y}_k=\{y_k,y_k(0),y_k(1),\lambda_k y_k,\lambda_k 


y_k(0),\lambda_k y_k(1)\}$. 


При этом ряд (11) сходится в норме пространства $W^1_2\times\Bbb 


C^2\times L_2\times\Bbb C^2$. 


Значит, тем более этот ряд сходится и в норме пространства $\frak H$.  


Коэффициенты $\{c_k\}$ в (10) и нужно  


было определить. Так как ряд из равенства (11) сходится в $\frak H$ к 


элементу  


$\wh{\bold y}$ и выполнены соотношения (8), то коэффициенты  


Фурье его разложения определяются однозначно и согласно  


построению элемента $\wh{\bold y}$ совпадают с коэффициентами $c_k$ из  


(10). Следовательно, для этих коэффициентов $c_k$


справедливa формулa (9). Теперь вычислим эти коэффициенты 


$(G\wh{\bold y},\wh{\bold z}_k)_{\frak H}=(v_0,\ -(\ov p_1z_k)'+ 


\ov q_1 z_k+\mu_k\ov q_2z_k)+


v_0(0)(a^{-1}_1\ov z_k(0)-p_1(0)-p_1(0)\ov z_k(0))+ 


v_0(1)(b^{-1}_1\ov z_k(1)-p_1(1)\ov z_k(1))+ 


(v_1,\ \ov q_2z_k)$. 


Воспользовавшись краевыми условиями сопряженной краевой задачи и  


тем, что $\mu_k=\ov\lambda_k$, получим равенства 


$(G\wh{\bold y},\wh{\bold z}_k)_{\frak H}=q^{-1}_2P_k$,


$(G\wh{\bold y}_k,\wh{\bold z}_k)_{\frak H}=q^{-1}_2Q_k$, 


откуда $c_k=P_k/Q_k$. 


 


Заметим, что второй  


способ вычисления коэффициентов является менее трудоемким. 


 


Автор благодарен профессору А.А.\,Шкаликову за постановку  


задачи и полезные обсуждения.  


 


\begin{thebibliography}{9} 


 


\bibitem{1} 


Будак Б.М., Самарский А.А., Тихонов А.Н. {\em Сборник задач по  


математической физике}. -- М.: Наука, 1980. -- 688 с.  


\bibitem{2} 


Шкаликов А.А. {\em Краевые задачи для обыкновенных  


дифференциальных уравнений с параметром в граничных условиях} // 


Тр. семин. им.~И.Г.\,Петровского. -- 1983. -- Вып.\,9. -- 


С.\,190--229.  


\bibitem{3} 


Ахтямов А.М. {\em О вычислении коэффициентов разложений по  


производным цепочкам одной спектральной задачи} // Матем. 


заметки. -- 1992. -- T.\,51. -- \No\,6. -- C.\,137--139.  


\bibitem{4} 


Шкаликов А.А. {\em Эллиптические уравнения в гильбертовом  


пространстве и спектральные задачи, связанные с ними} // Тр. 


семин. им.~И.Г.\,Петровского. -- 1989. -- Вып.\,14. -- С.\,140--224.  


\bibitem{5} 


Наймарк М.А. {\em Линейные дифференциальные операторы}. -- М.: 


Наука, 1967. -- 526 с. 


\bibitem{6} 


Михайлов В.П. {\em Дифференциальные уравнения в частных  


производных}. -- М.: Наука, 1976. -- 392 с.  


  


\end{thebibliography} 


 


\vskip 0.5cm 


 


\post{\it Башкирский государственный}{\it Поступила} 


\post{\it университет}{21.02.1996} 


 


\label{end} 


\end{document} 





